\chapter{The Derivation of the Vacuum Lattice Damping Rate ($\Gamma_D$)}
\label{app:vacuum_damping}
\addtocontents{toc}{\protect\setcounter{tocdepth}{0}}

\textbf{Subject:} Replacing the empirical Hubble Constant ($H_0$) with the theoretical Vacuum Damping Rate ($\Gamma_D$).

\section{The Physical Premise}
Standard cosmology attributes the redshift of distant galaxies to a recession velocity (metric expansion). Tension Field Theory (TFT) resolves this through the static vacuum's visco-elastic properties.

In TFT, the vacuum is a dual-phase Face-Centered Cubic (FCC) medium. Like all real visco-elastic media, it possesses intrinsic boundary friction.
\begin{itemize}
    \item \textbf{Mechanism:} A photon is a 2D transverse shear wave propagating exclusively across the solid FCC lattice nodes. \textit{(Note: This mechanism applies identically to 3D bulk gravitational waves)}. During propagation, the transverse wave must continuously work against the visco-elastic friction between these solid nodes and the interstitial condensate.
    \item \textbf{Result:} This work is not perfectly elastic. A fundamental, minute fraction of the wave's energy is dissipated into the lattice as heat (thermal exhaust). This boundary friction produces a strict temporal energy decay, which we define as the Vacuum Damping Rate ($\Gamma_D$). This cumulative energy loss is observed macroscopically as cosmological redshift ($z$).
\end{itemize}

\section{The Energy Dissipation Equation}
We define the energy loss rate as proportional to the total energy of the wave (a standard visco-elastic condensate drag characteristic):
\begin{equation}
    \frac{dE}{dt} = -\Gamma_D \cdot E
\end{equation}

\textbf{Where:}
\begin{itemize}
    \item $E$ is the photon energy ($h\nu$).
    \item $\Gamma_D$ is the Vacuum Damping Rate (units: $\text{s}^{-1}$).
\end{itemize}

Solving for a photon traveling over time $t$ yields the exponential decay:
\begin{equation}
    E(t) = E_0 e^{-\Gamma_D t}
\end{equation}

\section{Recovering the Linear Law}
For cosmic distances ($r$) where the total energy loss is small ($\Gamma_D t \ll 1$), we approximate the exponential decay as linear:
\begin{equation}
    \frac{\Delta E}{E} \approx -\Gamma_D \cdot t = -\Gamma_D \left(\frac{r}{c_E}\right)
\end{equation}

Since redshift is defined as $z = -\Delta E / E$ (fractional loss), we substitute:
\begin{equation}
    z = \left(\frac{\Gamma_D}{c_E}\right)r
\end{equation}

Rearranging to solve for the velocity-equivalent ($c_E z$):
\begin{equation}
    c_E \cdot z = \Gamma_D \cdot r
\end{equation}

\begin{itemize}
    \item \textbf{Observation:} This equation is mathematically identical to the Hubble Law ($v = H_0 d$), but the physical meaning is strictly inverted. $\Gamma_D$ is not a velocity of spatial expansion; it is a lattice-dynamic temporal damping rate.
\end{itemize}

\section{Calculation of $\Gamma_D$ from the Saturation Limit}
We derive the absolute magnitude of $\Gamma_D$ from the Saturation Condition. The Cosmic Horizon ($R_U$) is the physical distance where the cumulative vacuum drag exhausts the photon's energy budget.

\begin{itemize}
    \item \textbf{Condition:} At $r = R_U$, the integrated drag implies global topological saturation ($\alpha_{G,U} \to 1$).
    \item \textbf{Calculation:}
\end{itemize}
\begin{equation}
    \Gamma_D = \frac{c_E}{R_U}
\end{equation}

Using the exact optical velocity ($c_E = 299792458 \text{ m/s}$) and the high-precision observable causal radius locked by the duality pivot ($R_U \approx 1.2674 \times 10^{26} \text{ m}$):
\begin{equation}
    \Gamma_D = \frac{299792458}{1.2674 \times 10^{26}} \approx 2.3654 \times 10^{-18} \text{ s}^{-1}
\end{equation}

\section{Conclusion}
The Vacuum Damping Rate ($\Gamma_D$) is numerically equivalent to the empirical Hubble Constant ($H_0$).
\begin{equation}
    \Gamma_D \equiv H_0 \approx 2.3654 \times 10^{-18} \text{ s}^{-1}
\end{equation}
\textit{(Note: Converting this temporal rate into standard astronomical units yields exactly $72.987 \text{ km/s/Mpc}$, perfectly aligning with local-universe observations).}

\begin{itemize}
    \item \textbf{The Nature of Expansion:} The observed ``expansion'' of the universe is an artifact caused by the visco-elastic phase lag of the dual-phase spacetime lattice.
    \item \textbf{Universal Wave Drag:} Because this visco-elastic drag ($\Gamma_D$) is strictly a mechanical consequence of transverse geometric shear rubbing against the interstitial condensate, it applies identically to all transverse vacuum waves.
    \item \textbf{Gravitational Wave Resolution:} This provides the exact mechanical reason why macroscopic gravitational waves (3D bulk shear) suffer the exact same cosmological redshift as photons (2D surface shear) over cosmic distances.
    \item \textbf{LIGO Alignment:} This universal redshift rate ($z = (\Gamma_D/c_E)r$) perfectly aligns with the observational corrections required by the LIGO collaboration, proving the waves are bleeding energy into the visco-elastic SGF rather than being stretched by expanding space.
\end{itemize}